\section{Bounded block rank}
\label{sec:a-block-rank}

The cactus argument extends to networks with interacting cycles when
each biconnected block has bounded cycle rank; the number of blocks and
the total cycle rank may grow with the input. The extension has two
steps: localize free nominations to one block, then show that a small
objective perturbation leaves only a bounded number of free coordinates
inside it. Independent resistance intervals are then eliminated by a
parametric linear-programming argument in fixed dimension.

We write $\mathfrak r$ for the fixed block-rank bound, retaining
$\totalrank$ for total cycle rank. A selected block $H$ has entrance $a$
and exit $c_H$; the subscript distinguishes the exit vertex from the
weighted-objective vector $c$. Other local notation is defined where used.

\subsection{Results and scope}
Throughout, the graph is connected, the laws are symmetric quadratic,
and $\nomset=\{b:\ell\le b\le u,\ \one^\top b=0\}$ is a nonempty
balanced rational box. There are no additional operating constraints
on scenarios. All intervals are finite; singleton intervals are allowed.

\begin{theorem}[Bounded block rank with nomination uncertainty]
\label{thm:a-blk-nomination}
Fix an integer $\mathfrak r\ge0$. If $\blockrank\le\mathfrak r$,
resistances are positive rationals, and $s\ne t$, then MPD admits a
certified rational value interval of width at most $\eps$ and an
exactly feasible rational $\eps$-optimal nomination, in time polynomial
in $N$ and $\log(1/\eps)$ for fixed $\mathfrak r$.
The exponent may depend on $\mathfrak r$; this is an XP-type guarantee.
\end{theorem}

\begin{theorem}[Joint nomination and resistance intervals]
\label{thm:a-blk-joint}
The same conclusion holds jointly over
$\nomset\times\resset$, where
$\resset=\prod_{e\in E}[L_e,U_e]$ has rational $0<L_e\le U_e$.
The output scenario consists of rational nomination and resistance
vectors satisfying their bounds and nomination balance exactly.
\end{theorem}

\begin{theorem}[Exact arc extrema and capacity validation]
\label{thm:a-blk-arc}
Under the joint hypotheses of \cref{thm:a-blk-joint}, the maximum and
minimum signed flow on any prescribed arc are real algebraic numbers
of polynomial encoding length, computable exactly in polynomial bit
time for fixed $\mathfrak r$. Thus robust validation of all rational
signed arc capacities is decidable exactly in polynomial time for
fixed $\mathfrak r$, including equality at a capacity.
\end{theorem}

The joint uncertainty model is not new:
\citet[Section 3.3, equations (3.2)--(3.4), p.~56]{amann2019-exact-methods-for-two-stage}
uses the product of positive coefficient intervals and a balanced
nomination box. The resistance-only model and the circulation
sign-region construction appear in
\citet[Sections 4.1.4 and 4.3.2, Proposition 4.7
(arXiv v1 numbering)]{amann2018-deciding-robust-feasibility-and-infeasibility}.
The hypotheses here bound only the maximum rank of a block, allowing
unbounded total cycle rank and arbitrarily many uncertain coordinates.
The pressure and arc conclusions differ for arithmetic reasons: a
terminal drop can add values from many blocks, while an arc belongs to
one block and admits one algebraic representation of polynomial size.
Pressure is therefore handled by summing certified intervals, consistent
with the cactus SRS barrier. The exponent depends on the rank bound;
the theorems give tractability at each fixed rank, not a practically
small exponent.

\subsection{One active block and suppressed paths}
\begin{lemma}[Block ordering and localization]
\label{lem:a-blk-order}
Let the resistance vector be positive and real.
The aggregation statement of \cref{lem:a-cac-aggregation} holds on
any connected graph. On the resulting terminal core, the ordinary
smoothed adjoint has strictly ordered block ranges: every interior
vertex of a nonbridge block has value strictly between its entrance
and exit values. There is an $O(n)$ family of closed rational faces
containing an MPD maximizer such that each face either fixes the
nomination or leaves free coordinates in just one selected block.
Vertices strictly before that block are upper-saturated and those
strictly after it are lower-saturated. This family is independent of
resistances.
\end{lemma}
\begin{proof}
The aggregation proof, including its rational greedy disaggregation,
uses only the incidence tree and state uniqueness, so it applies verbatim.
The implicit-function and adjoint calculations in
\cref{lem:a-cac-adjoint} use only connectivity. The linear passive
state with nomination $e_s-e_t$ restricts in each path block to the
unit entrance--exit state, by \cref{lem:a-pre-block}.
Its entrance value exceeds its exit value: their difference equals
$\sum_{e\in E(H)}r_e\jmath_e^2>0$, since a unit flow is nonzero.
The weighted-average maximum principle bounds every other value by
the two boundary values. For strictness, suppose an interior vertex
has the entrance value. The block with the entrance deleted is
connected, so a path from this vertex to the exit avoids the entrance.
The weighted-average equality at successive interior vertices propagates
the maximum along this path to the exit, a contradiction. Deleting
the exit proves the strict minimum statement. A two-vertex block has
no interior vertices; a bridge has positive drop directly.
Hence open block ranges are disjoint and successive closed
ranges meet only at their shared articulation.

At a smoothed maximizer, feasible exchanges give the multiplier and
saturation rule \cref{eq:a-cac-saturation}; its proof allows fixed
coordinates and singleton polytopes. If the multiplier is inside a
block range, only vertices of that block can be free. At an articulation
or terminal value only that vertex can be free. A gap inside a bridge
leaves none. Above or below all ranges gives respectively the all-lower
or all-upper vector. Enumerate one face for each selected block, leaving
all its coordinates unrestricted and fixing vertices strictly outside
as stated, together with these vertex, gap, and extreme patterns.
Extra points in these faces are admissible and harmless; empty faces
are discarded by rational balance and interval tests.
The compactness and uniform smoothing argument in
\cref{lem:a-cac-faces} works on every connected graph. A subsequence of
smoothed maximizers stays in one of the finitely many closed faces;
its limit is an unsmoothed maximizer in that face.
\end{proof}

Fix a selected block. All other core nominations are fixed, and their
sums shift the two terminal intervals of the block to give the balanced
local box of aggregated nominations $b^H$. The sum of the block's
original core nominations is fixed by balance. Hence every other
block's aggregated nomination is fixed, even as the individual
nominations of this block vary, and for fixed resistances its
contribution is a constant by \cref{lem:a-pre-block}. The local
objective is $\pi_a-\pi_{c_H}$.

\begin{lemma}[Suppression topology]
\label{lem:a-blk-suppression}
For a nonbridge block put
$K_H=\{v:\deg_Hv\ge3\}\cup\{a,c_H\}$. Its edges partition into
$p_H$ maximal paths whose endpoints are in $K_H$ and whose internal
vertices have degree two. The interiors partition $V(H)\setminus K_H$,
and
\[
 |K_H|\le2r_H,\qquad p_H=r_H+|K_H|-1\le3r_H-1.
\]
\end{lemma}
\begin{proof}
The degree count \cref{lem:a-pre-degree} bounds the number of
high-degree vertices by $2r_H-2$; the two terminals give the first
bound. Every component after removing $K_H$ has maximum degree two.
It cannot be a cycle, since it would have no edge leaving it in the
connected block. Thus it is a path, possibly a singleton, with its two
missing incidences attached to $K_H$, and following a degree-two
vertex ends at marked vertices in both directions. They are distinct: a path
returning to one marked vertex would make that vertex a cut vertex,
since another marked terminal exists. Include also direct marked-to-marked
edges as paths without internal vertices. This partitions all edges.
Suppressing each interior removes equal numbers of edges and vertices,
preserves connectedness and cycle rank, and leaves a multigraph with
$|K_H|$ vertices and $p_H$ edges. This proves the equality and bound.
Parallel paths and rank-one blocks are included; bridges are excluded.
\end{proof}

\begin{lemma}[Positive forcing and unimodality]
\label{lem:a-blk-unimodal}
Use smoothed laws and perturb the local objective to
\[
 F_{\delta,\rho}(b^H)=\pi_a-\pi_{c_H}
       +\delta\sum_{v\in V(H)\setminus K_H}(\pi_v-\pi_{c_H}),\qquad \delta,\rho>0.
\]
On each suppressed path, its adjoint values first strictly increase
and then strictly decrease, allowing either part to be empty and a
maximum plateau of at most two adjacent vertices. Every level occurs
at most twice, and each strict upper-level set is an interval in path
order.
\end{lemma}
\begin{proof}
Linearity of the differentiated objective and the grounded calculation
of \cref{lem:a-cac-adjoint} give
\[
 L_\rho h=e_a-e_{c_H}
       +\delta\sum_{v\in V(H)\setminus K_H}(e_v-e_{c_H}),\qquad h_{c_H}=0.
\]
Orient a path as $v_0,\ldots,v_j$. If $R_i>0$ is the derivative
resistance on $v_iv_{i+1}$, set
$\jmath_i=(h_{v_i}-h_{v_{i+1}})/R_i$. Conservation at an internal
vertex $v_i$ says exactly
\[
 \jmath_i-\jmath_{i-1}=\delta>0,
\]
so the currents strictly increase. A negative current means an
increase in $h$, a positive current a decrease, and at most one current
is zero. A zero current separates these phases and gives a two-vertex
maximum plateau. Strictness on either side proves both level claims,
also for the subsequence of internal vertices and for coincident
endpoint values.
\end{proof}

\begin{lemma}[Bounded nomination faces]
\label{lem:a-blk-faces}
There is a resistance-independent family $\mathcal F$ of
$n^{O(\mathfrak r)}$ closed rational nomination faces containing an MPD
maximizer for every positive real resistance vector. In a selected
nonbridge block, each face has at most
$f_H\le |K_H|+2p_H\le8r_H-2$ free coordinates. Balance leaves at most
$8r_H-3$ independent nominations if $f_H>0$; if $f_H=0$ it only tests
feasibility. All other block nominations are fixed after aggregation.
Bridges and faces with a sole free coordinate can be reduced to fixed
rational nominations.
\end{lemma}
\begin{proof}
Refine each local box from \cref{lem:a-blk-order} separately. At a
maximizer of $F_{\delta,\rho}$ the same feasible-exchange argument gives
$h_v>\mu_0\Rightarrow b^H_v=u^H_v$ and
$h_v<\mu_0\Rightarrow b^H_v=\ell^H_v$.
Consequently the internal nominations of every path have the pattern
\[
 \text{lower / optional free pivot / upper /
 optional free pivot / lower}.
\]
Every segment may be empty. A two-vertex plateau at the threshold is
two adjacent free pivots with no intervening upper vertex.
For a path with $j$ internal vertices use positions $0,1,\ldots,2j$,
with vertices at odd positions and gaps at even positions. Choose
positions $i_1\le i_2$. Odd chosen positions are free, vertices strictly
between them upper, and the rest lower. These $O((j+1)^2)$ patterns
cover all thresholds, including all-upper, all-lower, and monotone
cases. Each marked vertex independently has status lower, upper, or
free. Enumerating all combinations gives at most
$3^{2r_H}(2n+1)^{2p_H}$ local patterns and the claimed free count.
Intersect with exact balance and remove empty patterns. Singleton
bounds are treated as fixed, regardless of the enumerated status.

The perturbed smoothed objectives converge uniformly. Orient each
nonzero physical flow by decreasing potential. It has no directed cycle;
successively subtracting the minimum flow on a source--sink path
produces a finite nonnegative path decomposition. Its total weight is
the total positive nomination. Thus $|x_e|\le B$, where on the aggregated
core we use the rational bound
\[
 B=\sum_v\max\{|\ell_v|,|u_v|\}.
\]
Block aggregation does not increase this bound. For $0<\rho\le1$,
normalized potentials are bounded in absolute value by
$(B^2+B)\sum_e\beta_e$, by summing edge drops on a simple path.
Therefore the perturbation is bounded by
$\delta n(B^2+B)\sum_e\beta_e$ uniformly on the local box.
The energy compactness argument of \cref{lem:a-cac-faces} gives uniform
convergence of the unperturbed smoothed objective. Letting both
parameters tend to zero and taking a subsequence in one closed pattern
face gives an original local maximizer in that face. Refining a coarse
face that already contains a global maximizer preserves its optimal
value because the other drops are constant. This is an existence proof;
the algorithm needs no perturbation size.

A bridge's local nomination is $(\zeta,-\zeta)$ with a rational feasible
interval for $\zeta$. Its entrance--exit drop $\beta\zeta|\zeta|$
is strictly increasing, so choose the largest endpoint. A sole free
coordinate is fixed by balance. Include these rational choices and
the fixed patterns from \cref{lem:a-blk-order}. Altogether the face
count is at most $C_{\mathfrak r}(n+1)^{6\mathfrak r+1}$ for a constant
$C_{\mathfrak r}$ when $\mathfrak r\ge1$; at rank zero it is $O(n)$.
\end{proof}

\subsection{Fixed resistances: algebraic local optimization}
On a face eliminate one free nomination by balance, if any. Let $z$
consist of the remaining nominations and the circulations in a
fundamental cycle basis of the selected block. Tree routing is zero
on every chord, and the basis has the identity matrix on chord rows.
Consequently
\begin{equation}
 x_e=X_e(z),\qquad
 d_H=\max\{0,f_H-1\}+r_H\le9r_H-3\quad(r_H\ge1),
 \label{eq:a-blk-core}
\end{equation}
where $X_e$ is rational affine, and each circulation is a chord flow
in $[-B,B]$. The nomination inequalities and this circulation box
form a compact rational polytope. An inactive block has only $r_H$
circulation variables. The dimension bounds for nonbridge blocks are
\[
\begin{array}{c|c}
\text{Quantity} & \text{Bound or definition}\\ \hline
|K_H| & \le 2r_H\\
p_H & \le 3r_H-1\\
f_H & \le 8r_H-2\\
d_H & \le 9r_H-3\\
d_* & =\max\{1,9\mathfrak r-3\}
\end{array}
\]

The circulation-coordinate cell count at fixed total rank is given by
\citet[Proposition 4.7 (arXiv v1 numbering)]{amann2018-deciding-robust-feasibility-and-infeasibility}.
Here the coordinates also include free nominations, and the rank bound
is imposed separately on each block.

The affine hyperplanes $X_e=0$ have at most
$\sum_{i=0}^{d_H}\binom{m_H}{i}$ full-dimensional cells, by the
standard recurrence. They can be constructed by inserting hyperplanes
successively and testing the two strict sides of each current cell by
rational linear feasibility; identically zero forms are omitted.
Thereafter use closed cell inequalities. Their union covers all of
space, even if the feasible polytope is lower dimensional or lies on a
zero-flow hyperplane. On a cell choose signs
$\epsilon_e\in\{-1,1\}$ and put $Q_e=\epsilon_e X_e^2$ (zero forms contribute
zero). These formulas agree on shared boundaries.

\begin{lemma}[Affine encoding on a sign cell]
\label{lem:a-blk-affine}
For fixed positive $\beta$, a core point $z$ satisfies the cell and cycle
constraints if and only if $X(z)$ is its nomination's physical flow
with those cell signs.
\end{lemma}
\begin{proof}
Write $Z_{H,ie}$ for the signed cycle incidence.
Let $\omega_{H,e}\in\{-1,0,1\}$ describe a fixed path from $a$ to $c_H$. The complete
physical and objective equations on a cell are
\begin{equation}
 \sum_e Z_{H,ie}\beta_e Q_e(z)=0\quad(1\le i\le r_H),
 \qquad G_H(z,\beta)=\sum_e\omega_{H,e}\beta_e Q_e(z).
 \label{eq:a-blk-polynomials}
\end{equation}
Indeed conservation is built into $X$, and orthogonality of the drop
vector to a basis of $\ker A_H$ is equivalent to its membership in
$\operatorname{im}A_H^\top$. Integrating along the tree constructs the
potentials, and uniqueness identifies these equations with the
physical state. They have degree two at fixed resistances.
\end{proof}

\begin{proof}[Proof of \cref{thm:a-blk-nomination}]
Aggregate the terminal core. Enumerate the faces in \cref{lem:a-blk-faces}.
At a rational threshold, \cref{eq:a-blk-polynomials} and the cell and
nomination constraints are a fixed-dimensional existential-real
system. Renegar's consistent-sign algorithm in
\cref{subsec:a-pre-computation} decides it in
$\tau^2(\sigma D)^{O(d_H+1)}$ bit operations, where here $D=2$,
$\sigma=O(N)$, and clearing denominators gives the safe bound
$\tau=O(N^4+\iota)$ at bisection step $\iota$.
Binary search with a retained feasible lower threshold gives an
active-block value interval of any requested width. Use the rational
absolute drop bound $\Vmax=B^2\sum_e\beta_e$ for the initial bracket.

For each inactive block the same equations have only $r_H$ variables
and a unique physical solution. Projection onto a circulation or drop
coordinate, using \cref{eq:a-pre-qe-cost}, and univariate isolation
therefore give exact algebraic coordinates and drops, with degrees
bounded in terms of $\mathfrak r$ and coefficient bits polynomial
in $N$. This also works for a solution on several sign boundaries:
project each closed cell, retain its unique real physical solution,
and discard empty cells. No nonsingular Jacobian or complex isolated-root
hypothesis is needed. A fixed bridge contributes a rational constant.
Isolate and refine block drops separately; never form the algebraic
field generated by all blocks.

Give each block an interval of width $\eps/(8n)$; summing gives each
face a certified interval of width at most $\eps/8$, and taking
separately the largest lower and upper endpoints gives a global
interval of that width. Choose a face with largest lower endpoint. Its retained active
threshold, together with the inactive block lower endpoints, is at
least that face lower endpoint. Recover an algebraic feasible active
core at this threshold by the fixed-dimension sample-point primitive
in \cref{prim:a-pre-sample-points}, applied to the cell, cycle, nomination,
and threshold constraints. It supplies all $d_H$ coordinates in one
common local algebraic representation, with polynomial degree and
encoding length for fixed $\mathfrak r$.

Put $C=2B\sum_{e\in P_{st}}\beta_e$ for any simple terminal path.
If $B=0$, all core flows and the objective vanish; greedily disaggregate
zero and finish. Otherwise refine each free nomination to a rational
isolating interval of width at most $\eps/(8nC)$. Intersect these
intervals with the face box and exact balance, keeping its fixed
coordinates. This rational polytope is nonempty because it contains
the algebraic nomination. Rational linear feasibility, or greedy
allocation between its coordinate bounds, produces a rational point.
Its $\ell_1$ error is at most $\eps/(8C)$, so
\cref{lem:a-cac-lipschitz} bounds the objective loss by $\eps/8$.
Choosing the face lost at most $\eps/8$; disaggregation loses nothing.
This proves both output guarantees, including fixed and singleton faces.

For explicit parameter-dependent bounds, set
$d_* =\max\{1,9\mathfrak r-3\}$ and
$N_\eps=N+\lceil\log_2(1/\eps)\rceil+1$, with $0<\eps\le1$.
There are at most $C_{\mathfrak r}(N+1)^{6\mathfrak r+1}$ faces,
at most $(N+1)^{d_*}$ cells per local problem, and at most $N$ blocks
per face. Bisection takes $O(N_\eps)$ steps. For an absolute constant
$A_0$, the total decision work is bounded by
\[
 C_{\mathfrak r}(N+1)^{6\mathfrak r+2+d_*+A_0(d_*+1)}
 (N^4+N_\eps)^2N_\eps.
\]
Projection of $d_H$ variables with one free value coordinate has
polynomial-count exponent $2(d_H+1)$ in \cref{eq:a-pre-qe-cost}.
A fixed number of such projections, sample-point calls, and root refinements
bounds all remaining work by
$C_{\mathfrak r}(N+N_\eps+1)^{2^{A_1(\mathfrak r+1)^2}}$
for an absolute constant $A_1$. These loose bounds include rational
arithmetic, output length, and rank-zero cases.
\end{proof}

\subsection{A parametric box lemma}
\begin{lemma}[Parametric box LP]
\label{lem:a-blk-boxlp}
Fix the core dimension $d_0$, the number $k$ of equations, and an input
degree bound $D_{\rm in}$. Let $z$ range over a compact rational semialgebraic
set $\mathcal C\subset\R^{d_0}$ described by an explicit quantifier-free
formula whose polynomials have degree at most $D_{\rm in}$. Let $\alpha_e(z),c_0(z)\in\Q[z]^k$ and
$g_e(z),g_0(z)\in\Q[z]$ have degree at most $D_{\rm in}$, and let
$\beta_e\in[L_e,U_e]$ ($1\le e\le m_0$) be independent finite rational intervals.
The attainable set
\[
 \mathcal V=\{(z,v):z\in\mathcal C,\ \exists\beta:
    \sum_e\beta_e \alpha_e(z)=c_0(z),\quad
    g_0(z)+\sum_e\beta_e g_e(z)=v\}
\]
has a quantifier-free description computable in polynomial bit time.
Feasibility and exact optimization of $v$ over this set, including an
algebraic optimizing core and leaf vector of polynomial encoding length,
are computable in polynomial bit time in the total rational input size.
\end{lemma}
\begin{proof}
Put $W_e(z)=(\alpha_e(z);g_e(z))$ and
$w(z,v)=(c_0(z);v-g_0(z))$. For fixed $z$ the image
\[
 \mathcal Z(z)=\sum_e W_e(z)[L_e,U_e]
\]
is compact, as the continuous image of the compact product box under
the sum map; convex, because coordinatewise convex combinations stay in
the box and the map is linear; and nonempty, also with zero columns and
singleton intervals. Its support function is
\[
 \max_{y\in\mathcal Z(z)}\eta^\top y
   =\sum_e\max\{L_e\eta^\top W_e(z),U_e\eta^\top W_e(z)\}.
\]
The equality follows by independently choosing a maximizing endpoint
of each interval; ties allow either endpoint.
A point $w$ belongs to a nonempty compact convex set $S$ if and only
if it satisfies every support inequality. Necessity is immediate;
for sufficiency, if $w\notin S$, choose a closest point $y_0\in S$.
For each $y\in S$, differentiation of the squared distance along
$y_0+\vartheta(y-y_0)$ at $\vartheta=0+$ gives
$(w-y_0)^\top(y-y_0)\le0$. Taking $\eta=w-y_0$ gives
$\eta^\top w=\eta^\top y_0+\|w-y_0\|_2^2>
\max_{y\in S}\eta^\top y$, a contradiction. Hence membership is exactly
\begin{equation}
 z\in\mathcal C\quad\wedge\quad
 \forall\eta\in\R^{k+1}:\quad
 \eta^\top w(z,v)\le
 \sum_e\max\{L_e\eta^\top W_e(z),U_e\eta^\top W_e(z)\}.
 \label{eq:a-blk-support}
\end{equation}

We expand the maxima without introducing one variable per maximum or
listing all formal sign assignments. Enumerate only the
realizable sign vectors of the polynomials
$t_e(z,\eta)=\eta^\top W_e(z)$ in $d_0+k+1$ variables, using Renegar's
bound in \cref{subsec:a-pre-computation}. For each realized vector
choose $U_e$ for positive sign and $L_e$ for negative or zero sign.
On its sign condition replace the right side by the resulting sum,
and take the disjunction of these guarded inequalities. This is
exact at zero as well. With $m_0$ leaves there are
$(m_0(D_{\rm in}+1))^{O(d_0+k+1)}$ such vectors, polynomially many terms,
and degree at most $\max\{2,D_{\rm in}+1\}$. Summation and rational denominator clearing
have polynomial coefficient bit length. This supplies a polynomial-size
first-order formula with $d_0+1+(k+1)$ variables.
\Cref{eq:a-pre-qe-cost}, with one quantified block of size $k+1$ and
$d_0+1$ free variables, has polynomial-count exponent
$(k+2)(d_0+2)$ and fixed degree and bit factors. It gives the asserted
quantifier-free description.

The full feasible set is closed in the product of the compact core
and the boxes, hence compact, and so is its image $\mathcal V$.
If $\mathcal V$ is nonempty, the maximum is attained. Append to the quantifier-free
formula $\Psi(z,v)$ the assertion
\[
 \neg\exists z',v':\ \Psi(z',v')\ \wedge\ v'>v.
\]
This again has a fixed number of variables. Eliminate the quantifiers
using \cref{eq:a-pre-qe-cost}, then apply the sample-point primitive
in \cref{prim:a-pre-sample-points}
to obtain an exact optimal pair $(z^*,v^*)$ in the common real algebraic
field $\Q(\theta_{\rm alg})$, with polynomial degree and encoding length.
Empty sets are detected first.

The leaves are recovered constructively, so that no exponential
enumeration of box corners is hidden in an LP step. All $W_e(z^*)$
and $w(z^*,v^*)$ lie in that same ordered algebraic field.
Enumerate the full-dimensional cells of the central hyperplanes
$\eta^\top W_e(z^*)=0$ for nonzero columns with $L_e<U_e$.
There are $m_0^{O(k+1)}$ cells. Fixed-dimensional exact linear sign tests
can be done in $\Q(\theta_{\rm alg})$, or by adjoining its single defining
root and isolating inequalities to the rational tests. For each cell
record the maximizing endpoint vector $\beta^{(j)}$ and its image
$y_j=\sum_e W_e(z^*)\beta_e^{(j)}$. Zero columns and singleton boxes
have a fixed arbitrary endpoint choice.
These images include every vertex of $\mathcal Z(z^*)$. A singleton
zonotope is immediate: every support choice has its sole image. Otherwise,
the image of a box is the convex hull of the images of its finitely
many corners, hence a polytope. A vertex can be strictly separated
from the convex hull of the other distinct corner images. Its exposing
directions contain an open set; choosing one outside the finitely many
nonzero hyperplanes yields a cell in the enumeration. Conversely its
cell support choice gives that vertex. Thus their convex hull is the
whole zonotope, also when its affine dimension is smaller than $k+1$.

Every point in a convex hull in $\R^{k+1}$ is a convex combination of
at most $k+2$ affinely independent vertices. Indeed, in a representation
with a dependent support, an affine dependence lets one change the
weights in both directions while preserving their sum and the point;
move until one weight vanishes and repeat. Enumerate subsets of at
most $k+2$ of the $y_j$, solve their affine systems, and test nonnegative
weights. At least one passes for $w(z^*,v^*)$. For its weights $\varpi_j$
return $\beta=\sum_j\varpi_j\beta^{(j)}$, which lies in the box and has
exactly the required image, so attains $v^*$. Only fixed-size linear
systems and polynomially many field operations are used. Exact sign
tests and this fixed-size linear algebra take polynomial bit time in
$\Q(\theta_{\rm alg})$, using the root's isolating representation and
rational polynomial arithmetic. The solutions are determinant ratios of
fixed size, and the final sums have polynomial encoding length in the
same field. Ties, redundant
equations, zero columns, and singleton zonotopes require no genericity.

An optimal vertex of the original box LP with $k$ equations has at
most $k$ non-bound coordinates: otherwise a dependence of the
corresponding columns permits a two-sided feasible motion. This fact
alone is not a polynomial algorithm, because choosing those coordinates
still leaves exponentially many endpoint assignments. The recovery
above avoids that enumeration: it uses only the polynomially many
exposed zonotope vertices and need not output an LP vertex.
\end{proof}

\subsection{Joint optimization and rational recovery}
\begin{lemma}[Joint Lipschitz estimate]
\label{lem:a-blk-resistance-lipschitz}
On any connected graph with the stated boxes, put
$B=\sum_v\max\{|\ell_v|,|u_v|\}$ and fix a simple terminal path $P_{st}$.
For all admissible scenarios in these unfiltered product boxes,
\[
 |F_{s,t}(b,\beta)-F_{s,t}(b',\beta')|
 \le C_b\|b-b'\|_1+B^2\|\beta-\beta'\|_1,
 \qquad C_b=2B\sum_{e\in P_{st}}U_e.
\]
\end{lemma}
\begin{proof}
The nomination estimate \cref{lem:a-cac-lipschitz} is uniform after
replacing each resistance by its upper bound. For resistance sensitivity
hold $b$ fixed and differentiate the smoothed equations. With
$R_\rho=\operatorname{diag}(\beta_e(2|x_e|+\rho))$ and
$T_e=x_e|x_e|+\rho x_e$, they give
\[
 A\,dx=0,\qquad A^\top d\pi=R_\rho\,dx+
                       \operatorname{diag}(T_e)\,d\beta.
\]
Multiply by $AR_\rho^{-1}$, then pair with the ordinary unit adjoint
$L_\rho h=e_s-e_t$. Symmetry yields
\[
 dF_\rho=\sum_e\jmath_e T_e\,d\beta_e,
 \qquad \frac{\partial F_\rho}{\partial\beta_e}
          =\jmath_e(x_e|x_e|+\rho x_e).
\]
There is no sign reversal, since the parameter term is on the right of
the differentiated potential equation. The nonzero electrical currents
follow strictly decreasing $h$, so they are acyclic, and conservation
has just one unit source and one unit sink. Following positive edges from the
source to the sink and subtracting the minimum edge current repeatedly
decomposes them into paths of total weight one. Therefore
$|\jmath_e|\le1$. The physical path decomposition in
\cref{lem:a-blk-faces} gives $|x_e|\le B$, whence the derivative is
bounded by $B^2+\rho B$.
Integrate on the resistance-box segment. Uniform convergence as
$\rho\downarrow0$ holds jointly over both compact boxes: in the energy
argument use the common positive lower bound $\min_e L_e$ and the
upper bounds $U_e$ for the tree comparison energies. Uniqueness
identifies every subsequential limit as before, and passing to the
limit gives the resistance estimate. Changing nominations first and
resistances second, the triangle inequality completes the proof.
\end{proof}

\begin{proof}[Proof of \cref{thm:a-blk-joint}]
A joint maximizer exists by \cref{lem:a-pre-attainment}. With its
resistances fixed, \cref{lem:a-blk-faces} supplies an equally good
nomination in the resistance-independent family, which is jointly
optimal, since an improvement would contradict the original joint
maximum.
On any selected face, inactive aggregated nominations are fixed.
\Cref{lem:a-pre-block} and independence of resistance intervals make
the face maximum the sum of independently optimized block drops.

For an active nonbridge cell use \cref{eq:a-blk-core}. Apply
\cref{lem:a-blk-boxlp} with $k=r_H$, $\alpha_{e,i}=Z_{H,ie}Q_e$,
$c_0=0$, $g_e=\omega_{H,e} Q_e$, and $g_0=0$. The core is the compact
rational nomination/circulation polytope intersected with the closed
sign cell. All coefficient degrees are at most two. The cycle equations
are exactly the linking equations, so both directions of the mapping
follow from \cref{eq:a-blk-polynomials}. An inactive block uses only
its $r_H$ circulations and the same lemma. For a bridge, write its
flow in entrance--exit orientation as $\zeta$. At fixed $\zeta$ the
best resistance is $U_e$ if $\zeta\ge0$ and $L_e$ if $\zeta<0$.
The best value remains strictly increasing in $\zeta$, so an active
bridge takes the largest feasible $\zeta$; zero flow permits either
endpoint. This covers fixed flows and singleton intervals.

Optimize each cell exactly, compare its algebraic values within that
block, and retain a local optimizer. Approximate each block optimum by
an interval of width $\eps/(8n)$, sum by face, and select the largest
face lower endpoint. The resulting global interval has width at most
$\eps/8$, and the selected face has an exact optimizing scenario with
loss at most $\eps/8$. Each block's algebraic data are stored
separately; sums are never compared and fields never combined.

If $B=0$ on the aggregated core, its objective is zero: disaggregate
zero nominations and choose rational resistance endpoints everywhere
(off-core flows need not vanish). Otherwise $C_b>0$. Refine every free
core nomination coordinate to an interval of width
$\eps/(8nC_b)$ and recover a rational balanced point in its face by
the rational box/balance feasibility procedure above. Refine each
relevant resistance coordinate to width $\eps/(8mB^2)$, intersect with
$[L_e,U_e]$, and choose any rational point there. These operations give
nomination loss at most $\eps/8$ and resistance loss at most $\eps/8$
by \cref{lem:a-blk-resistance-lipschitz}. The total loss, including
selection, is at most $3\eps/8$. Greedy disaggregation and arbitrary
rational off-core resistance choices complete the output. Circulations
need not survive rounding: the rounded scenario has its own physical
state, and the Lipschitz estimate compares the objectives.

For bit costs, the face, block, and cell counts above still apply.
In the box lemma, $d_0\le d_*$ and $k\le\mathfrak r$, so the sign
construction has $N^{O(d_*+\mathfrak r+1)}$ terms and its first QE
has count exponent $(\mathfrak r+2)(d_*+2)$. Its output degrees depend
only on $\mathfrak r$. The optimal-pair formula adds one fixed-size
quantifier block; the sample-point primitive in
\cref{prim:a-pre-sample-points} then supplies the
optimal core and value in fixed dimension. These sampling bounds and
\cref{eq:a-pre-qe-cost} bound the full algorithm,
including algebraic recovery and refinement, by
\[
 C_{\mathfrak r}(N+N_\eps+1)^{2^{A_2(\mathfrak r+1)^2}}
\]
for an absolute constant $A_2$, which is polynomial for every fixed
rank bound. All constants used in rounding have polynomial rational bit
length; the factors $n,m$ cost only logarithmic extra precision.
\end{proof}

\subsection{Exact local arc values}
\begin{proof}[Proof of \cref{thm:a-blk-arc}]
Let $e=(u,v)$ belong to block $H$. The block path between its endpoints
consists of $H$ alone. Aggregate all outside components into their
attachments. The resulting balanced interval box has exactly the
attainable local nominations, and outside resistances have no effect
on $x_e$, by \cref{lem:a-pre-block,lem:a-cac-aggregation}.
At fixed resistances,
\[
 x_e=\sgn(\pi_u-\pi_v)\sqrt{|\pi_u-\pi_v|/\beta_e}
\]
is strictly increasing in its endpoint drop. Fixing the resistances of
a joint flow maximizer therefore transfers a maximizer to the faces of
\cref{lem:a-blk-faces} for terminals $(u,v)$.
On each cell the target flow $X_e(z)$ is affine. In
\cref{lem:a-blk-boxlp} use the cycle coefficients as before, but now
$g_0=X_e$ and all $g_f=0$. This gives an exact algebraic cell optimum
of polynomial encoding length. Compare the polynomially many cell
and face values exactly by real-algebraic sign determination.
To minimize, reverse the ordered terminals and the target orientation
and maximize $-X_e$ with the corresponding face family.

For a bridge let $S_e$ be its tail component and $\bar S_e=V\setminus S_e$.
Writing $\ell(S_e)=\sum_{v\in S_e}\ell_v$, and similarly for the other
sums, the exact attainable flow interval is
\[
 [\ell(S_e),u(S_e)]\cap[-u(\bar S_e),-\ell(\bar S_e)].
\]
Interval addition and greedy allocation prove attainability of every
point, and nonemptiness follows from the balanced-box hypothesis. The
rational endpoints are independent of resistances.

Repeat for every tested edge and check whether its maximum is at most
its rational upper capacity and its minimum at least its lower capacity;
exact sign tests include equality. There are at most $2m$ optimizations
with the same fixed-rank polynomial bound. Capacities test all original
scenarios and never filter the optimization domain.
Unlike the pressure span in \cref{thm:a-cac-srs}, an arc objective
involves no sum of independently encoded block values: all its
algebraic computation occurs in one bounded-dimensional block core.
\end{proof}

\begin{remark}[Optional rational arc witnesses]
\label{rem:a-blk-arc-witness}
Round an algebraic \emph{arc-flow} optimizer, with target flow $x$,
and let $y$ be the target flow of the rounded scenario. If their endpoint
drops are $\Delta,\Delta'$ and target resistances $\beta_e,\beta'_e$,
then
\[
 |x-y|^2\le
 \frac{2}{L_e}\bigl(|\Delta-\Delta'|+B^2|\beta_e-\beta'_e|\bigr).
\]
Indeed $|x|x-|y|y$ has magnitude at least $|x-y|^2/2$: for equal signs
use $|x+y|\ge|x-y|$, and for opposite signs use
$x^2+y^2\ge(|x|+|y|)^2/2$.
Subtracting $\Delta=\beta_ex|x|$ and
$\Delta'=\beta'_ey|y|$ gives
$\beta_e|x|x-\beta_ey|y|=\Delta-\Delta'+(\beta'_e-\beta_e)y|y|$;
take absolute values and use $|y|\le B$ and $\beta_e\ge L_e$.
The joint Lipschitz estimate for $(u,v)$ then permits rational scenario
rounding to a desired flow error, using its square in the error budget
and only polynomially many precision bits.
A pressure-maximizing resistance scenario cannot replace the arc
optimizer: on two parallel edges with unit total flow, other resistance
one, and target resistance in $[1,4]$, the target flow is
$1/(1+\sqrt{\beta_e})$, decreasing in $\beta_e$, whereas its endpoint
drop $\beta_e/(1+\sqrt{\beta_e})^2$ is increasing.
\end{remark}

\subsection{Limits of the structural and arithmetic claims}
\begin{proposition}[Broader graph parameters do not suffice]
\label{prop:a-blk-limits}
Unless $\mathrm P=\NP$, a polynomial additive algorithm with polynomial
dependence on accuracy bits cannot hold on all networks of treewidth
two, on all series-parallel networks, or on all networks of feedback
vertex number one, even with fixed positive rational resistances.
\end{proposition}
\begin{proof}
Use \citet[Section 4, Figure 2, equation (3), and Lemmas 4.3 and 4.17
(numbering of the April 6, 2022 manuscript)]{thurauf2022-deciding-the-feasibility-of-a}.
For a Partition instance with $j\ge3$ positive integers $S_i$ and
$K_0=\sum_iS_i$, his graph has edges
$sz_i^+,z_i^+t,z_i^+z_i^-$.
The bookings are $K_0/2$ at $s,t$ and $S_i$ at $z_i^+,z_i^-$;
translate entries to intervals $[0,b_v]$ and exits to $[-b_v,0]$.
The main-edge resistances are $1/S_i^2$ and pendant resistances one.
These are admissible rational data here, with a nonempty balanced box,
and his objective (2) imposes no potential bounds. The cited lemmas
give maximum drop at least one for a yes instance and strictly below
$\Theta_0(K_0)<1$ for a no instance. For an explicit encoding gap,
equation (3) defines, in renamed notation,
\begin{align*}
 \vartheta_0&=1-\left(1-\frac{1}{8K_0^2}\right)^2,&
 \Pi_0&=\max\{1-\vartheta_0+\vartheta_0^2,\ 1-\vartheta_0^2/K_0^2\},\\
 \widetilde\vartheta_0&=(1-\Pi_0)/5,&
 \Theta_0(K_0)&=\max\{1-\widetilde\vartheta_0^2/(K_0^2j^2),\ \Pi_0+4\widetilde\vartheta_0\}.
\end{align*}
Here $0<\vartheta_0<1$, $0<\Pi_0<1$, and both terms defining $\Theta_0(K_0)$ are
strictly below one. Fixed powers, rational arithmetic, and maxima give
polynomial encoding length, so $\mathfrak g=1-\Theta_0(K_0)>0$ is rational with
polynomially bounded $\log(1/\mathfrak g)$. A certified interval of width $\mathfrak g/4$
lies entirely above $(1+\Theta_0(K_0))/2$ in the yes case and below it in the
no case, deciding Partition.

By \cref{lem:a-pre-ranks}, the main block $K_{2,j}$ has rank $j-1$
and is series-parallel; adding pendant edges preserves that property. Bags
$\{s,t,z_i^+\}$ in a path, with each bag $\{z_i^+,z_i^-\}$ attached
to its corresponding main bag, give a width-two tree decomposition.
There is a cycle, so the treewidth is exactly two. Deleting $s$ leaves
the tree with branches $t-z_i^+-z_i^-$, and the cycle shows that the
feedback vertex number is exactly one. These class bounds and the gap
prove the proposition. This does not contradict the bounded-block-rank
results, since $j-1$ grows.
\end{proof}

\begin{remark}[Convex fibers and exact comparison]
\label{rem:a-blk-convex-barrier}
Convexity alone cannot replace the box hypothesis in the exact-output
lemma without encountering the SRS barrier. With no nonlinear core,
let leaf $i$ be the singleton
$\{y_i:0\le y_i\le a_i+1,\ y_i^2=a_i\}$ for a positive integer $a_i$.
Each fiber is compact and convex and has a degree-two rational
description, but its unique point is $\sqrt{a_i}$. Exact optimization
and rational-threshold comparison of $\sum_i y_i$ would decide
$\sum_i\sqrt{a_i}\le K$, and these independent fibers supply no fixed
common algebraic field of bounded degree. Thus the box lemma's proof
does not extend merely by assuming convex quadratic fibers. This is an
exact-arithmetic obstruction, not an NP-hardness claim or a proof that
every larger leaf class is intractable.
\end{remark}
